\chapter{Proposed Methodology}
\label{ch:methodology}

This chapter describes the end-to-end system. Section~\ref{sec:overview} gives the overall design, Section~\ref{sec:formulation} formalizes the task, and Sections~\ref{sec:data}--\ref{sec:degradation} cover the data, the log-mel frontend, and the degradation simulator. Sections~\ref{sec:baseline}--\ref{sec:star-edge-method} present the baseline, the proposed FT-RCL detector, and the STAR-Edge extension. Sections~\ref{sec:optimization}--\ref{sec:evaluation-metrics} describe optimization, the difference between training and inference, and the evaluation metrics.

\section{System overview}
\label{sec:overview}

The proposed system keeps a fixed, compact backbone (MobileNetV3-Small~\cite{howard2019mobilenetv3}) and improves its behaviour under channel changes only through the training objective. Figure~\ref{fig:pipeline-overview} summarizes the workflow: waveforms are segmented, converted to normalized log-mel features, and classified; during robustness training each waveform is additionally degraded and passed through the same network, and the two views are tied together by coherence losses.

\subsection{Design principles}

\begin{itemize}
  \item \textbf{Unchanged inference graph.} All robustness mechanisms act at training time; at test time the model computes one mel spectrogram and runs one forward pass.
  \item \textbf{Clean baseline for attribution.} A clean-trained baseline (mode~M0, class-weighted cross-entropy on clean log-mel inputs only) uses identical preprocessing, data, and architecture, so that any effect of the auxiliary losses can be attributed to them.
  \item \textbf{Consistency learning.} Two views of the same instance, a clean waveform and a degraded variant with the same label, should map to similar outputs. Penalizing disagreement between the views counteracts unstable predictions and unstable internal features, whereas augmentation alone only exposes the model to degraded data without requiring agreement.
\end{itemize}

% Master proposed-methodology flowchart (replaces pipeline_overview.png)
\begin{figure}[htbp]
\centering
\resizebox{\textwidth}{!}{%
\begin{tikzpicture}
% (a) data preparation and paired views
\node[blk,text width=7.6cm] (in) at (0,0) {\textbf{AT-ADD Track~2 clip}\\mono, 24\,kHz, label $y\in\{0,1\}$, audio type $t$};
\node[blk,text width=9.2cm] (seg) at (0,-1.6) {\textbf{Segmentation}\\4\,s crop (random in training, centre in evaluation) or repeat-padding, $N=96{,}000$};
\node[blk,text width=4.6cm] (xc) at (-3.8,-3.3) {\textbf{Clean view} $x$};
\node[trn,text width=6.4cm] (g) at (3.7,-3.3) {\textbf{Degradation simulator} $\tilde{x}=g(x;c)$\\MP3 $\cdot$ noise $\cdot$ resampling $\cdot$ telephony;\\severity $\rho\le\rho_{\max}(e)$ (curriculum)};
% (b) frontend and backbone
\node[blk,text width=5.6cm] (mc) at (-3.8,-6.3) {\textbf{Log-mel frontend} $\phi$\\STFT $\to$ 80 mel bins $\to$ log $\to$ z-score\\output $\hat{\mathbf L}_c$};
\node[trn,text width=5.6cm] (md) at (3.8,-6.3) {\textbf{Log-mel frontend} $\phi$ (identical)\\STFT $\to$ 80 mel bins $\to$ log $\to$ z-score\\output $\hat{\mathbf L}_d$};
\node[blk,text width=10.4cm] (bb) at (0,-8.6) {\textbf{Shared MobileNetV3-Small backbone} ($\approx$1.52\,M parameters)\\inverted residual blocks $\to$ feature map $F$ $\to$ global pool $\to$ 2-class head};
\node[blk,text width=5.2cm,minimum height=1.5cm] (oc) at (-3.8,-10.7) {\textbf{Clean outputs}\\logits $z_c$, score $q_c$,\\8-band signature $s_c$};
\node[trn,text width=5.2cm,minimum height=1.5cm] (od) at (3.8,-10.7) {\textbf{Degraded outputs}\\logits $z_d$, score $q_d$,\\8-band signature $s_d$};
% (c) objectives
\node[los,text width=3.2cm,minimum height=1.9cm] (lcls) at (-5.8,-14.0) {$\mathcal{L}_{\mathrm{cls}}$\\dual weighted CE\\$\mathcal{L}_{\mathrm{CE}}(z_c,y)+\beta\mathcal{L}_{\mathrm{CE}}(z_d,y)$};
\node[los,text width=3.2cm,minimum height=1.9cm] (lpc) at (-1.95,-14.0) {$\mathcal{L}_{\mathrm{PC}}$\\prediction coherence\\$(q_c-q_d)^2$};
\node[los,text width=3.2cm,minimum height=1.9cm] (lft) at (1.95,-14.0) {$\mathcal{L}_{\mathrm{FTSC}}$\\signature coherence\\$1-\cos(s_c,s_d)$ per band};
\node[los,text width=3.2cm,minimum height=1.9cm] (gate) at (5.8,-14.0) {\textbf{Robustness gate}\\$w_i=\alpha_{t_i}\kappa_i\eta(\rho_i)$\\type, confidence, severity};
\node[los,text width=12.6cm] (tot) at (0,-16.5) {\textbf{Total objective (mode M3)}\\$\mathcal{L}_{\mathrm{M3}}=\mathcal{L}_{\mathrm{cls}}+\bar{w}\,(\gamma_{\mathrm{pc}}\mathcal{L}_{\mathrm{PC}}+\gamma_{\mathrm{ftsc}}\mathcal{L}_{\mathrm{FTSC}})$};
% (d) optimisation and (e) inference side by side
\node[blk,text width=6.8cm] (opt) at (-3.9,-19.6) {\textbf{Optimization}\\AdamW, lr $10^{-3}$, weight decay $10^{-4}$, batch 32, $\le$20 epochs; gradients update the shared backbone};
\node[blk,text width=6.8cm] (sel) at (-3.9,-21.6) {\textbf{Robustness-aware checkpoint selection}\\$J=\mathrm{Macro\text{-}F1}_{\mathrm{dev}}-\lambda\,\widehat{\Delta\mathrm{F1}}_{\mathrm{dev}}$, $\lambda=0.5$};
\node[dep,text width=6.8cm] (inf) at (3.9,-19.6) {\textbf{Deployed detector} (inference graph unchanged)\\one log-mel pass $\to$ one forward pass\\$q=p_1$, $\hat y=\mathbb 1[q\ge0.5]$};
\node[dep,text width=6.8cm] (ev) at (3.9,-21.6) {\textbf{Evaluation} on held-out clips\\Macro-F1, accuracy, AUROC, EER,\\AT-ADD Macro-F1};
% arrows
\draw[arr] (in)--(seg);
\draw[arr] ([xshift=-1.6cm]seg.south)--(xc.north);
\draw[arr] ([xshift=1.6cm]seg.south)--(g.north);
\draw[arr] (xc)--(mc);
\draw[arr] (g)--(md);
\draw[arr] (mc.south)--(mc.south |- bb.north);
\draw[arr] (md.south)--(md.south |- bb.north);
\draw[arr] (oc.north)--(oc.north |- bb.south);
\draw[arr] (od.north)--(od.north |- bb.south);
\draw[thick] (oc.south)--(oc.south |- 0,-12.1);
\draw[thick] (od.south)--(od.south |- 0,-12.1);
\draw[thick] (-5.8,-12.1)--(5.8,-12.1);
\foreach \n in {lcls,lpc,lft,gate}{\draw[arr] (\n.north |- 0,-12.1)--(\n.north);}
\foreach \n in {lcls,lpc,lft,gate}{\draw[arr] (\n.south)--(\n.south |- tot.north);}
\draw[arr] (tot.south -| opt.north)--(opt.north);
\draw[arr] (opt)--(sel);
\draw[arr] (sel.east)--++(0.45,0)|-node[pos=0.25,left,font=\scriptsize,align=right]{best\\checkpoint}(inf.west);
\draw[arr] (inf)--(ev);
% stage frames
\begin{scope}[on background layer]
\node[stage,fit=(in)(seg)(xc)(g)] (sA) {};
\node[stage,fit=(mc)(md)(bb)(oc)(od)] (sB) {};
\node[stage,fit=(lcls)(lpc)(lft)(gate)(tot)] (sC) {};
\node[stage,fit=(opt)(sel)] (sD) {};
\node[stage,fit=(inf)(ev),draw=deployedge] (sE) {};
\end{scope}
\node[stlab,fill=white,anchor=center] at (sA.north) {(a) Data preparation and paired views};
\node[stlab,fill=white,anchor=center] at (sB.north) {(b) Shared frontend and backbone};
\node[stlab,fill=white,anchor=center] at ([xshift=-4.6cm]sC.north) {(c) FT-RCL training objectives};
\node[stlab,fill=white,anchor=west] at ([xshift=0.3cm]sD.north west) {(d) Optimization};
\node[stlab,fill=white,anchor=center] at (sE.north) {(e) Inference and evaluation};
\end{tikzpicture}}
\caption{Overall end-to-end workflow of the proposed methodology. Orange blocks are used only during training; the deployed detector (green) keeps the single-view baseline inference graph.}
\label{fig:pipeline-overview}
\end{figure}


\section{Problem formulation}
\label{sec:formulation}

Let $x \in \mathbb{R}^{N}$ be a mono waveform with sampling rate $f_s = 24{,}000$\,Hz and fixed length $N = 4 f_s = 96{,}000$ samples (four seconds). The label $y \in \{0,1\}$ indicates real ($0$) or fake ($1$). The classifier produces logits $z \in \mathbb{R}^2$ and probabilities $p = \mathrm{softmax}(z)$. The \textbf{detection score} used for ROC analysis and EER is the fake-class probability $q = p_1$.

For robustness training, a condition $c$ from a finite catalogue defines a waveform-domain operator $g(\cdot\,; c)$. The paired sample $(x, \tilde{x})$ satisfies $\tilde{x} = g(x; c)$ with the same label $y$ and the same temporal crop. Mel extraction $\phi(\cdot)$ is applied \textbf{identically} to both waveforms:
\begin{equation}
\mathbf{X} = \phi(x), \qquad \tilde{\mathbf{X}} = \phi(\tilde{x}).
\label{eq:pair}
\end{equation}
Degradation never alters the mel hyperparameters or normalization statistics, which ensures a fair comparison with the clean baseline.

\section{Dataset}
\label{sec:data}

\subsection{AT-ADD Track~2}

Experiments use \textbf{AT-ADD Track~2}, which provides binary real/fake labels with metadata for the audio type (speech, singing, music, sound) and, for fake clips, the synthesis generator. The benchmark defines separate train, development, and evaluation partitions, and the fake generators in evaluation are disjoint from those in training~\cite{xie2026atadd}. Real and fake clips are imbalanced in training (roughly eight real clips for every fifteen fake clips), which is handled by class weighting (Section~\ref{sec:baseline}).

\subsection{Experimental data}

Table~\ref{tab:dataset-eval} summarizes the size of the AT-ADD Track~2 benchmark, the held-out evaluation clips used for reported metrics, and the main training settings. The manifests contain 126{,}825 clips in total. Type counts in the evaluation subset are not balanced, so pooled metrics are dominated by speech.

\begin{table}[htbp]
\centering
\caption{Dataset size, evaluation subset, and training settings.}
\label{tab:dataset-eval}
\begin{tabular}{@{}llr@{}}
\toprule
\textbf{Group} & \textbf{Quantity} & \textbf{Value} \\
\midrule
\multirow{4}{*}{Full benchmark (clips)}
  & Training split (16{,}000 real, 30{,}000 fake) & 46{,}000 \\
  & Development split & 36{,}536 \\
  & Evaluation split & 44{,}289 \\
  & \textbf{Total} & \textbf{126{,}825} \\
\midrule
\multirow{5}{*}{Eval subset by audio type (clips)}
  & Speech & 8{,}083 \\
  & Singing & 2{,}023 \\
  & Sound & 1{,}794 \\
  & Music & 1{,}387 \\
  & \textbf{Total evaluated} & \textbf{13{,}287} \\
\midrule
\multirow{2}{*}{Training settings}
  & Batch size & 32 \\
  & Random seed & 42 \\
\bottomrule
\end{tabular}
\end{table}

\section{Audio preprocessing and log-mel frontend}
\label{sec:frontend}

\subsection{Loading and segmentation}

Clips are resampled to $f_s$ if necessary. If the duration exceeds four seconds, a crop of length $N$ is taken (uniform random start in training, centre crop in evaluation). Shorter clips are extended by \textbf{repeat-padding} (tiling the waveform) rather than zero-padding.

\subsection{Short-time Fourier transform and mel filtering}

With a Hann window $w[n]$, frame length $L_w = 600$ samples (25\,ms), hop $H = 240$ samples (10\,ms), and FFT size $N_{\mathrm{FFT}} = 1024$ (the window is zero-padded to the FFT size),
\begin{equation}
X(m,k) = \sum_{n=0}^{L_w-1} x[n + mH]\, w[n]\, e^{-j2\pi kn / N_{\mathrm{FFT}}}, \qquad P(m,k) = |X(m,k)|^2 .
\label{eq:stft}
\end{equation}
Frames are centred using reflection padding, so a four-second segment yields $T_m = 1 + \lfloor N/H \rfloor = 401$ frames. A mel filter bank $\mathbf{W}_{\mathrm{mel}} \in \mathbb{R}^{80 \times 513}$ with 80 bands (Slaney normalization) gives the mel power $\mathbf{P}_{\mathrm{mel}} = \mathbf{W}_{\mathrm{mel}} \mathbf{P}^{\top}$, and the log-mel features are
\begin{equation}
\mathbf{L} = \log(\mathbf{P}_{\mathrm{mel}} + \varepsilon), \qquad \varepsilon = 10^{-10}.
\label{eq:logmel}
\end{equation}

\subsection{Global normalization}

The mean $\mu$ and standard deviation $\sigma$ are computed on the training split ($\mu \approx -9.19$, $\sigma \approx 6.85$ in the reference baseline) and applied in training and testing:
\begin{equation}
\hat{\mathbf{L}} = (\mathbf{L} - \mu) / \sigma .
\label{eq:norm}
\end{equation}
The tensor $\hat{\mathbf{L}}$ is treated as a single-channel image of shape $1 \times 80 \times T_m$ by the CNN. Figure~\ref{fig:mel-frontend} illustrates the frontend.

\begin{figure}[htbp]
\centering
\resizebox{\linewidth}{!}{%
\begin{tikzpicture}
\node[blk,text width=2.1cm] (w) at (0,0) {\textbf{Waveform}};
\node[blk,text width=1.8cm] (s) at (2.5,0) {\textbf{STFT}};
\node[blk,text width=2.0cm] (m) at (5.0,0) {\textbf{Mel bank}};
\node[blk,text width=1.5cm] (l) at (7.5,0) {\textbf{log}};
\node[blk,text width=2.0cm] (z) at (10.0,0) {\textbf{Z-score}};
\draw[arr] (w)--(s);
\draw[arr] (s)--(m);
\draw[arr] (m)--(l);
\draw[arr] (l)--(z);
\end{tikzpicture}}
\caption{Audio preprocessing and log-mel feature extraction.}
\label{fig:mel-frontend}
\end{figure}


\section{Synthetic degradation (paired views)}
\label{sec:degradation}

\subsection{Degradation catalogue}

Degradations are applied in the waveform domain before $\phi$. Table~\ref{tab:degradations} lists the training catalogue; each condition has a severity level $\rho$ that increases with the strength of the degradation. Figure~\ref{fig:spectrogram-pairs} shows a clean and a degraded version of one clip in the waveform and log-mel domains. MP3 encoding requires an external codec tool on the host system. Each degradation records the condition family, severity level, and random seed, which are used by the curriculum and the gate.

\begin{table}[htbp]
\centering
\caption{Degradation families and severity levels (training catalogue).}
\label{tab:degradations}
\small
\begin{tabular}{@{}lll@{}}
\toprule
Family & Setting & Severity $\rho$ \\
\midrule
Resample & 24\,kHz $\rightarrow$ 16\,kHz $\rightarrow$ 24\,kHz & 1 \\
MP3 codec & 128 / 64 / 32\,kbps & 2 / 3 / 4 \\
Additive noise & SNR 20 / 10 / 5\,dB & 2 / 3 / 4 \\
Telephony & Band-pass 300--3400\,Hz & 3 \\
\bottomrule
\end{tabular}
\end{table}

\figimage{spectrogram-pairs}{Clean versus degraded waveform and log-mel on the evaluation split (MP3 degradation is shown only when the codec tool is available).}{../results/figures/btp_report/spectrogram_pairs.png}

\subsection{Severity curriculum}

Let $\rho_{\max}(e)$ be the maximum severity allowed at epoch $e$, with milestones at epochs 6 and 13:
\begin{equation}
\rho_{\max}(e) = \begin{cases}
1 & e < 6, \\
3 & 6 \le e < 13, \\
4 & e \ge 13 .
\end{cases}
\label{eq:curriculum}
\end{equation}
Only conditions with $\rho \le \rho_{\max}(e)$ are sampled, so early training sees only the mildest degradation and the harder conditions of Table~\ref{tab:degradations} enter gradually (Figure~\ref{fig:m3-gating}).

\figimage{m3-gating}{Severity curriculum of mode M3: maximum allowed severity $\rho_{\max}$ versus epoch.}{../results/figures/btp_report/m3_gating_curriculum.png}

\subsection{Severity scaling}

For a sample with severity $\rho$, the severity factor used by the M3 gate (Section~\ref{sec:losses}) is
\begin{equation}
\eta(\rho) = \min\!\left(1,\ \frac{\rho}{\rho_{\max}(e)}\right).
\label{eq:eta}
\end{equation}

\section{Baseline model (mode M0)}
\label{sec:baseline}

\subsection{Architecture}

The baseline uses \textbf{MobileNetV3-Small} adapted to single-channel mel input (without ImageNet pre-training) and a two-logit classifier head, giving $1{,}519{,}618$ parameters. The forward pass is $\hat{\mathbf{L}} \mapsto z = f_\theta(\hat{\mathbf{L}})$ (Figure~\ref{fig:mobilenet-baseline}). Its depthwise-separable convolutions and inverted bottlenecks keep the inference cost low. As demonstrated empirically in Section~\ref{subsec:baseline-arch-ablation}, pure 2D convolutional modeling with spatial global average pooling outperforms recurrent CRNN architectures (LSTM/GRU) on multi-domain audio by over $2.2\%$ EER, establishing MobileNetV3-Small as the optimal lightweight backbone for all-type audio.

\begin{figure}[htbp]
\centering
\resizebox{\linewidth}{!}{%
\begin{tikzpicture}
\node[blk,text width=2.4cm] (in) at (0,0) {$1\times80\times T$\\log-mel};
\node[blk,text width=3.4cm] (mn) at (3.6,0) {\textbf{MobileNetV3-Small}\\inverted residuals};
\node[blk,text width=2.2cm] (gp) at (7.8,0) {\textbf{Global pool}};
\node[blk,text width=2.2cm] (hd) at (11.0,0) {\textbf{2-class head}};
\draw[arr] (in)--(mn);
\draw[arr] (mn)--(gp);
\draw[arr] (gp)--(hd);
\end{tikzpicture}}
\caption{MobileNetV3-Small baseline on single-channel log-mel input (${\sim}1.52$\,M parameters).}
\label{fig:mobilenet-baseline}
\end{figure}


\subsection{Class-weighted cross-entropy loss}

With class counts $n_0, n_1$ and total $n = n_0 + n_1$, the weights $\omega_j = n / (2 n_j)$ balance the real and fake classes:
\begin{equation}
\mathcal{L}_{\mathrm{CE}}(z, y) = -\sum_{j \in \{0,1\}} \omega_j \,\mathbb{1}[y=j]\,\log p_j .
\label{eq:ce}
\end{equation}
For the training class ratio of 16:30 (real:fake) this gives $\omega_0 \approx 1.44$ and $\omega_1 \approx 0.77$. Mode \textbf{M0} (clean-only training) minimizes $\mathcal{L}_{\mathrm{CE}}$ on clean batches only.

\section{FT-RCL detector}
\label{sec:ftrcl}

\subsection{Architecture and subband signatures}

The \textbf{FT-RCL detector} shares the MobileNet backbone and exposes both the classification logits and an intermediate feature map $F \in \mathbb{R}^{B \times C \times H \times W}$ (Figure~\ref{fig:ft-rcl-paired}). \textbf{SubbandPool} partitions the frequency axis $H$ of the feature map into $K=8$ bands. For band $k$ with rows $[h_k, h_{k+1})$,
\begin{equation}
b_k = \mathrm{mean}_{h,w}\big(F_{:,:,h_k:h_{k+1},:}\big) \in \mathbb{R}^{B \times C}.
\label{eq:subband}
\end{equation}
Stacking gives $s^{(0)} = [b_1, \ldots, b_K] \in \mathbb{R}^{B \times K \times C}$, which is L2-normalized along the channel axis:
\begin{equation}
s = \mathrm{normalize}\big(s^{(0)},\ p=2,\ \mathrm{dim}=-1\big).
\label{eq:signature}
\end{equation}
This gives a lightweight frequency-time summary of the internal representation, in contrast to full self-supervised embeddings.

\begin{figure}[htbp]
\centering
\resizebox{0.92\linewidth}{!}{%
\begin{tikzpicture}
\node[blk,text width=2.6cm] (c) at (-2.2,1.1) {\textbf{Clean mel}};
\node[trn,text width=2.6cm] (d) at (-2.2,-1.1) {\textbf{Degraded mel}};
\node[blk,text width=3.4cm] (bb) at (1.6,0) {\textbf{Shared MobileNetV3}\\backbone + signatures};
\node[los,text width=2.0cm] (lcls) at (5.4,1.35) {$\mathcal{L}_{\mathrm{cls}}$\\dual CE};
\node[los,text width=2.0cm] (lpc) at (5.4,0) {$\mathcal{L}_{\mathrm{PC}}$\\coherence};
\node[los,text width=2.0cm] (lft) at (5.4,-1.35) {$\mathcal{L}_{\mathrm{FTSC}}$\\signatures};
\draw[arr] (c)--(bb);
\draw[arr] (d)--(bb);
\draw[arr] (bb)--(lcls);
\draw[arr] (bb)--(lpc);
\draw[arr] (bb)--(lft);
\end{tikzpicture}}
\caption{FT-RCL paired clean/degraded-view training architecture.}
\label{fig:ft-rcl-paired}
\end{figure}


\subsection{Loss functions}
\label{sec:losses}

Let $z_c, z_d$ be the logits from the clean and degraded mel inputs, $q_c, q_d$ the corresponding fake-class probabilities, and $s_c, s_d$ the signatures.

\subsubsection{Dual classification}

With the class-weighted cross-entropy of Equation~\eqref{eq:ce},
\begin{equation}
\mathcal{L}_{\mathrm{cls}} = \mathcal{L}_{\mathrm{CE}}(z_c, y) + \beta\,\mathcal{L}_{\mathrm{CE}}(z_d, y), \qquad \beta = 0.5 .
\label{eq:cls}
\end{equation}

\subsubsection{Prediction coherence (PC)}

\begin{equation}
\mathcal{L}_{\mathrm{PC}} = \frac{1}{B}\sum_{i=1}^{B} \big(q_{c,i} - q_{d,i}\big)^2 .
\label{eq:pc}
\end{equation}
For example, suppose that for one clip the fake-class probabilities are $q_c = 0.80$ and $q_d = 0.55$. Then $(q_c - q_d)^2 = 0.0625$ contributes to $\mathcal{L}_{\mathrm{PC}}$, and the term pushes the degraded prediction toward the clean one. If the degraded view hides high-frequency artifacts, PC penalizes the resulting overconfident disagreement. PC does not by itself guarantee a correct decision, since both views could agree on a wrong label.

\subsubsection{Feature signature coherence (FTSC)}

For band $k$ and sample $i$, let $s_{c,i,k}, s_{d,i,k} \in \mathbb{R}^{C}$ be the L2-normalized band vectors and $\cos_{i,k} = s_{c,i,k}^\top s_{d,i,k}$. FTSC compares the clean and degraded signatures by the cosine similarity along the channel dimension for each band and averages over bands and batch:
\begin{equation}
\mathcal{L}_{\mathrm{FTSC}} = \frac{1}{B}\sum_{i=1}^{B} \frac{1}{K}\sum_{k=1}^{K} (1 - \cos_{i,k}).
\label{eq:ftsc}
\end{equation}
Figure~\ref{fig:pc-ftsc} illustrates PC and FTSC.

\figimage{pc-ftsc}{Prediction coherence and feature signature coherence.}{../results/figures/btp_report/pc_ftsc.png}

\subsubsection{Type, confidence, and severity robustness gate}

For sample $i$ with audio type $t_i$, the gate combines a type weight $\alpha_{t_i}$ (defaults: speech and singing $1.0$, music $0.3$, sound $0.5$), a confidence $\kappa_i = 2|q_{c,i} - 0.5|$ (treated as a constant, so no gradient flows through $\kappa$), and the severity factor $\eta(\rho_i)$ of Equation~\eqref{eq:eta}:
\begin{equation}
w_i = \alpha_{t_i}\,\kappa_i\,\eta(\rho_i), \qquad \bar{w} = \frac{1}{B}\sum_{i=1}^{B} w_i .
\label{eq:gate}
\end{equation}
The gate limits over-regularization on speech-dominated batches and lets the coherence terms act most strongly on confident, strongly degraded samples.

\subsection{Total loss and training modes}

The total loss of the full FT-RCL objective (mode M3) is
\begin{equation}
\mathcal{L}_{\mathrm{M3}} = \mathcal{L}_{\mathrm{cls}} + \bar{w}\,\Big(\gamma_{\mathrm{pc}}\mathcal{L}_{\mathrm{PC}} + \gamma_{\mathrm{ftsc}}\mathcal{L}_{\mathrm{FTSC}}\Big), \qquad \gamma_{\mathrm{pc}} = 1.0,\ \gamma_{\mathrm{ftsc}} = 0.5 .
\label{eq:m3}
\end{equation}
Table~\ref{tab:modes} summarizes how the four training modes combine the terms of Equations~\eqref{eq:cls}--\eqref{eq:ftsc} and the gate of Equation~\eqref{eq:gate}. Mode~M1 uses $\mathcal{L}_{\mathrm{CE}}(z_d,y)$ only; mode~M2 uses $\mathcal{L}_{\mathrm{cls}} + \bar{w}\,\gamma_{\mathrm{pc}}\mathcal{L}_{\mathrm{PC}}$ with $\eta \equiv 1$ and $\gamma_{\mathrm{ftsc}} = 0$.

\begin{table}[htbp]
\centering
\caption{Training modes and their loss components.}
\label{tab:modes}
\small
\renewcommand{\arraystretch}{1.2}
\begin{tabular}{@{}llll@{}}
\toprule
Mode & Network & Views used in the loss & Loss \\
\midrule
M0 & MobileNetV3-Small & Clean & $\mathcal{L}_{\mathrm{CE}}(z_c,y)$ \\
M1 & FT-RCL detector & Degraded & $\mathcal{L}_{\mathrm{CE}}(z_d,y)$ \\
M2 & FT-RCL detector & Clean and degraded & $\mathcal{L}_{\mathrm{cls}} + \bar{w}\,\gamma_{\mathrm{pc}}\mathcal{L}_{\mathrm{PC}}$ ($\eta\equiv 1$) \\
M3 & FT-RCL detector & Clean and degraded & $\mathcal{L}_{\mathrm{cls}} + \bar{w}\,(\gamma_{\mathrm{pc}}\mathcal{L}_{\mathrm{PC}} + \gamma_{\mathrm{ftsc}}\mathcal{L}_{\mathrm{FTSC}})$ \\
\bottomrule
\end{tabular}
\end{table}

\section{STAR-Edge extension}
\label{sec:star-edge-method}

STAR-Edge extends the FT-RCL recipe with three additional components. It uses the same backbone, frontend, degradation catalogue, and optimization settings, and is enabled through its own configuration file.

\subsection{Asymmetric bona-fide robustness (ABR)}

ABR applies an additional robustness term to \emph{real} clips. Genuine recordings are the class most affected when compression or noise introduces artifacts that resemble synthesis, so the term treats bona-fide clips asymmetrically to keep them from being flagged as fake after degradation.

\subsection{Severity trajectory coherence (STC)}

STC extends pairwise coherence across increasing severity levels of the same clip, so that the representation changes consistently as the degradation strengthens rather than only matching a single clean--degraded pair.

\subsection{Channel auxiliary detection (CAD)}

CAD derives an eight-dimensional channel code from waveform statistics and trains an auxiliary head to detect the channel condition. The backbone embeddings are projected orthogonally to a learned channel direction before signature extraction, which discourages the signatures from encoding channel information instead of evidence of synthesis. As with the other training-only components, CAD is not used at inference. The training progress of the STAR-Edge run is reported in Section~\ref{sec:star-edge}.

\section{Optimization and checkpoint selection}
\label{sec:optimization}

\subsection{Optimizer and training schedule}

Training uses AdamW with learning rate $10^{-3}$ and weight decay $10^{-4}$, batch size 32, up to 20 epochs, mixed-precision training on CUDA, and early stopping with patience 5 on a selection score. The learning rate follows ReduceLROnPlateau on development Macro-F1 (factor 0.5, patience 2).

\subsection{Checkpoint selection}

\textbf{Baseline.} The checkpoint with the best development Macro-F1 is kept.

\textbf{FT-RCL ($J$-score).} Let $\widehat{\Delta \mathrm{F1}}_{\mathrm{dev}}$ be an estimate of the Macro-F1 drop under quick on-the-fly degradations on a subset of development batches. The selection score is
\begin{equation}
J = \mathrm{MacroF1}_{\mathrm{dev,clean}} - \lambda\,\widehat{\Delta \mathrm{F1}}_{\mathrm{dev}}, \qquad \lambda = 0.5 ,
\label{eq:jscore}
\end{equation}
and checkpoints maximize $J$ when robustness-aware selection is enabled.

\section{Training versus inference}
\label{sec:train-infer}

\textbf{Training (M3)} requires paired degradations, two mel computations, dual logits and signatures, and the full loss $\mathcal{L}_{\mathrm{M3}}$. \textbf{Inference} applies $\phi$ once and performs a single forward pass; signatures and degradations are omitted. This preserves the latency expected of an edge-compatible deployment (Figure~\ref{fig:train-vs-infer}).

\begin{figure}[htbp]
\centering
\begin{tikzpicture}
\begin{scope}
\node[font=\scriptsize\bfseries\itshape,text=black!70,anchor=south west] at (-0.15,1.35) {Training (M3)};
\draw[dashed,rounded corners=4pt,draw=black!45] (-0.35,-0.85) rectangle (4.85,1.05);
\node[blk,text width=1.4cm] (tm) at (0,0.35) {\textbf{Mel}};
\node[blk,text width=1.4cm] (tc) at (2.2,0.35) {\textbf{CNN}};
\node[blk,text width=1.4cm] (ts) at (4.4,0.35) {\textbf{Loss}};
\draw[arr] (tm)--(tc);
\draw[arr] (tc)--(ts);
\node[font=\scriptsize\itshape,text=black!70,align=center,text width=4.6cm] at (2.2,-0.45) {Paired clean/degraded views\\dual CE + PC + FTSC + gate};
\end{scope}
\begin{scope}[xshift=6.5cm]
\node[font=\scriptsize\bfseries\itshape,text=black!70,anchor=south west] at (-0.15,1.35) {Inference (deploy)};
\draw[dashed,rounded corners=4pt,draw=deployedge] (-0.35,-0.85) rectangle (4.85,1.05);
\node[dep,text width=1.4cm] (im) at (0,0.35) {\textbf{Mel}};
\node[dep,text width=1.4cm] (ic) at (2.2,0.35) {\textbf{CNN}};
\node[dep,text width=1.4cm] (is) at (4.4,0.35) {\textbf{Score}};
\draw[arr] (im)--(ic);
\draw[arr] (ic)--(is);
\node[font=\scriptsize\itshape,text=black!70,align=center,text width=4.6cm] at (2.2,-0.45) {Single mel pass\\single forward pass};
\end{scope}
\end{tikzpicture}
\caption{Training versus inference workflow.}
\label{fig:train-vs-infer}
\end{figure}


\section{Evaluation metrics}
\label{sec:evaluation-metrics}

\subsection{Metrics}

For a clip with label $y \in \{0,1\}$, the model produces logits $z$ and probabilities $p = \mathrm{softmax}(z)$ (Section~\ref{sec:formulation}). The \textbf{detection score} is
\begin{equation}
q = p_1 = \frac{e^{z_1}}{e^{z_0} + e^{z_1}} \in [0,1],
\label{eq:q}
\end{equation}
the predicted probability that the clip is fake. For a decision threshold $\tau$, the predicted label is $\hat{y}_\tau = \mathbb{1}[q \ge \tau]$. Unless stated otherwise, $\tau = 0.5$. Let $\mathrm{TP}$, $\mathrm{FP}$, $\mathrm{TN}$, and $\mathrm{FN}$ denote the counts of true positives, false positives, true negatives, and false negatives, with the fake class ($y=1$) as the positive class. The metrics are defined as follows.

\paragraph{Accuracy:} the fraction of correct decisions at threshold $0.5$ on $q$,
\begin{equation}
\mathrm{Acc} = \frac{\mathrm{TP} + \mathrm{TN}}{\mathrm{TP} + \mathrm{TN} + \mathrm{FP} + \mathrm{FN}}.
\label{eq:acc}
\end{equation}

\paragraph{Macro-F1:} the unweighted mean of the F1 scores of the real and fake classes. For class $c \in \{0,1\}$, with precision $P_c$ and recall $R_c$,
\begin{equation}
\mathrm{F1}_c = \frac{2\,P_c R_c}{P_c + R_c}, \qquad
\mathrm{Macro\text{-}F1} = \tfrac{1}{2}\left(\mathrm{F1}_0 + \mathrm{F1}_1\right),
\label{eq:macrof1}
\end{equation}
where $P_1 = \mathrm{TP}/(\mathrm{TP}+\mathrm{FP})$, $R_1 = \mathrm{TP}/(\mathrm{TP}+\mathrm{FN})$, $P_0 = \mathrm{TN}/(\mathrm{TN}+\mathrm{FN})$, and $R_0 = \mathrm{TN}/(\mathrm{TN}+\mathrm{FP})$. Unless stated otherwise, the overall values are computed over all evaluation clips pooled together.

\paragraph{AUROC:} the area under the ROC curve traced by the score $q$. With the false-positive rate $\mathrm{FPR}(\tau) = \mathrm{FP}/(\mathrm{FP}+\mathrm{TN})$ and the true-positive rate $\mathrm{TPR}(\tau) = \mathrm{TP}/(\mathrm{TP}+\mathrm{FN})$ at threshold $\tau$,
\begin{equation}
\mathrm{AUROC} = \int_{0}^{1} \mathrm{TPR}\big(\mathrm{FPR}^{-1}(u)\big)\,du
= \Pr\big(q_{+} > q_{-}\big) + \tfrac{1}{2}\Pr\big(q_{+} = q_{-}\big),
\label{eq:auroc}
\end{equation}
where $q_{+}$ and $q_{-}$ are the scores of a randomly drawn fake and real clip, respectively.

\paragraph{EER:} the error rate at the operating point where the false-positive and false-negative rates are equal. With $\mathrm{FNR}(\tau) = \mathrm{FN}/(\mathrm{FN}+\mathrm{TP})$, it is computed at the ROC point where $\mathrm{FPR}$ and $\mathrm{FNR}$ are closest,
\begin{equation}
\tau^{\ast} = \arg\min_{\tau}\,\big|\mathrm{FPR}(\tau) - \mathrm{FNR}(\tau)\big|, \qquad
\mathrm{EER} = \tfrac{1}{2}\big(\mathrm{FPR}(\tau^{\ast}) + \mathrm{FNR}(\tau^{\ast})\big).
\label{eq:eer}
\end{equation}

\paragraph{AT-ADD Macro-F1:} for each audio type $t \in \mathcal{T} = \{\text{speech}, \text{singing}, \text{music}, \text{sound}\}$, the F1 scores of the real and fake classes are averaged, and the four type-wise scores are then averaged without weighting~\cite{xie2026atadd},
\begin{equation}
\mathrm{F1}_{\mathrm{AT\text{-}ADD}} = \frac{1}{|\mathcal{T}|}\sum_{t \in \mathcal{T}} \tfrac{1}{2}\left(\mathrm{F1}_{0}^{(t)} + \mathrm{F1}_{1}^{(t)}\right),
\label{eq:atadd}
\end{equation}
where $\mathrm{F1}_{c}^{(t)}$ is the F1 score of class $c$ computed on the clips of type $t$ only. Reported results use clean held-out evaluation clips only; no hyperparameters are tuned on evaluation data.
